% TOP_PROBLEM: {"id": 3184, "title": "Jorgensen's Conjecture", "category_id": 3, "category": {"id": 3, "name": "graph_theory", "display_name": "Graph Theory", "description": "Problems involving graphs, networks, and their properties.", "slug": "graph-theory", "order_index": 3, "created_at": "2026-07-31T15:26:25.671Z"}, "status": "open", "problem_number": "OPG-154", "set_id": 12}
% FIRST_POSTED: 2026-10-01
% ATTEMPT_STATUS: unresolved
% UnsolvedMath ID 3184; source catalogue code OPG-154
% !TeX program = lualatex
% GPT-6 Astra Ultra research attempt, 2026-10-01; independently checked partial work.
% Completion estimate: no numerical percentage assigned; see final status and literature audit.
\documentclass[11pt,a4paper]{article}
\usepackage[margin=25mm]{geometry}
\usepackage{fontspec}
\setmainfont{lmroman10-regular.otf}[BoldFont=lmroman10-bold.otf,
 ItalicFont=lmroman10-italic.otf,BoldItalicFont=lmroman10-bolditalic.otf]
\usepackage{amsmath,amssymb,mathtools}
\usepackage{unicode-math}
\setmathfont{latinmodern-math.otf}
\AtBeginDocument{\let\setminus\smallsetminus}
\usepackage{xurl,hyperref}
\hypersetup{colorlinks=true,urlcolor=blue}
\setcounter{secnumdepth}{0}
\setlength{\parindent}{0pt}
\setlength{\parskip}{5pt}
\setlength{\emergencystretch}{3em}
\begin{document}
\section*{Jorgensen's Conjecture}\label{3184}
{\small UnsolvedMath ID 3184 · OPG-154 · Graph Theory · OpenGarden}
\subsection{Definitions and mathematical statement}
All graphs in this statement are finite, simple, and
undirected. A graph $G$ is $6$-connected if it has at
least seven vertices and deleting any set of at most
five vertices leaves a connected graph. A minor of
$G$ is obtained by deleting vertices or edges and
contracting edges; loops and repeated edges created
by contraction may be removed.

Equivalently, $G$ has a $K_6$ minor when it has six
pairwise disjoint nonempty vertex sets
$B_1,\ldots,B_6$, each inducing a connected subgraph,
such that an edge joins $B_i$ to $B_j$ for every
distinct pair $i,j$. The sets are the branch sets
contracted to the six vertices of the complete minor.

A graph is apex if deleting at most one vertex makes
it planar. Planarity means the resulting graph can
be drawn in the plane with edges meeting only at
their common endpoints. Jørgensen's conjecture is
\[
 G\text{ $6$-connected and without a $K_6$ minor}
 \quad\Longrightarrow\quad
 \exists v\in V(G):G-v\text{ is planar}.
\]
No bound on $|V(G)|$ is included: every graph meeting
the two hypotheses is covered. The apex vertex may
depend on $G$, and no planar embedding or apex
candidate is provided in advance. Absence of a minor
is stronger than absence of a $K_6$ subgraph, since
contracting connected branch sets is permitted.
\subsection{Short English statement}
Does deleting one suitable vertex make every 6-connected graph without a complete six-vertex minor planar?
\subsection{Research attempt}
\paragraph{Scope and process.}
GPT-6 Astra Ultra, 1 October 2026. The catalogue formulation is retained above. This attempt checks the original problem and primary literature, proves the restricted claims stated below, and identifies the exact limit of its conclusions. Reproved facts are not claimed as novel results.

\paragraph{Literature and proof target.}
Kawarabayashi, Norine, Thomas and Wollan [S3] prove J\o rgensen's conjecture for all sufficiently large 6-connected graphs. That theorem is a substantial external result, not a proof supplied here of all finite orders. We audit the easy direction, obtain exact small-order exclusions, and derive restrictions that any apex vertex in a 6-connected graph must satisfy. Graphs in this attempt are finite and simple, as in the catalogue.

\paragraph{Why apex graphs have no $K_6$ minor.}
Suppose $G-v$ is planar, and assume there is a model of $K_6$ in $G$ with six pairwise disjoint connected branch sets. At most one branch set contains $v$. The other five are connected within $G-v$, remain pairwise adjacent there, and would give a $K_5$ minor in a planar graph. This is impossible: deletions and contractions preserve planarity, whereas $K_5$ has ten edges on five vertices, exceeding the planar bound $3\cdot5-6=9$.

This establishes apex $\Rightarrow K_6$-minor-free without connectivity hypotheses. The conjecture asks for a converse under 6-connectivity, so using this direction backwards would be a logical error. In particular, the absence of a $K_6$ model does not itself identify the desired vertex.

\paragraph{Restrictions on an apex in a 6-connected graph.}
If $G$ is 6-connected and $G-v$ is planar, deleting any further four vertices leaves $G-v$ connected. Thus $G-v$ is at least 5-connected and has minimum degree at least five. Put $m=|V(G-v)|$. Degree summation and the planar edge bound give
\[
 \frac{5m}{2}\le |E(G-v)|\le3m-6,
 \qquad m\ge12. \tag{1}
\]
Hence a 6-connected apex graph has at least thirteen vertices. More generally, any candidate apex must leave a planar graph with all these connectivity and degree restrictions; merely finding a planar-looking drawing after deleting several vertices is insufficient.

There is also a converse construction at the level of examples. Given any 5-connected planar graph $H$, add one universal vertex $v$. A deletion of at most five vertices leaves the result connected: if $v$ remains it connects all remaining vertices, and if $v$ is deleted then at most four vertices of $H$ were also deleted. This gives a 6-connected apex graph and, by the first paragraph, excludes a $K_6$ minor. Thus the two hypotheses are compatible; the size restriction is not a contradiction to the existence of the proposed class.

\paragraph{Direct verification at orders seven and eight.}
A 6-connected graph has at least seven vertices and minimum degree six. At order seven it is $K_7$, so contains $K_6$ as a subgraph. At order eight every vertex has at most one non-neighbour; its complement is a matching together with isolated vertices. Extend that missing-edge matching to a perfect matching on all eight vertices by arbitrarily pairing the remaining, evenly many vertices. Label its pairs $\{a_i,b_i\}$, $1\le i\le4$. Then $G$ contains the complete four-partite graph with these four pairs as parts.

Inside that subgraph take the six branch sets
\[
 \{a_1\},\{a_2\},\{a_3\},\{a_4\},
 \{b_1,b_2\},\{b_3,b_4\}.
\]
The two nonsingleton sets are connected. The four singleton sets are mutually adjacent. Every singleton has a neighbour in each doubleton, because a doubleton meets two different parts; the two doubletons are adjacent as well. These are therefore a $K_6$ model. This rules out $K_6$-minor-free 6-connected graphs at both orders, without invoking the general conjecture.

\paragraph{Verification, gap, and final status.}
The companion script \texttt{scripts/check\_attempt\_batch02\_graphs\_2026\_10\_01.py} checks the displayed branch sets in the eight-vertex complete four-partite graph. Additional edges can only preserve a minor certificate, which justifies all smaller missing-edge matchings. Equation (1) is a necessary condition on apex graphs; it does not prove that every 6-connected graph of orders nine through twelve contains a $K_6$ minor. Nor do the small certificates address the finite exceptions potentially left outside the sufficiently-large literature theorem. These are explicit gaps. The unrestricted catalogue implication is \textbf{UNRESOLVED} by this attempt.

\subsection{Sources}
\begin{itemize}
\item[S1] ULAM.AI and the UnsolvedMath Contributors, supplied problems.json, record 3184 (OPG-154), OpenGarden collection. Catalogue identity and source transcription; CC BY 4.0.
\url{https://huggingface.co/datasets/ulamai/UnsolvedMath}.
\item[S2] Open Problem Garden, Jorgensen's Conjecture. Original problem and discussion; the mathematical formulation below is expanded from the supplied source record.
\url{http://www.openproblemgarden.org/op/jorgensens_conjecture}.
\item[S3] Ken-ichi Kawarabayashi, Serguei Norine, Robin Thomas and Paul Wollan, \emph{$K_6$ minors in large 6-connected graphs}, arXiv:1203.2192 (2012). The stated theorem covers all sufficiently large graphs.
\url{https://arxiv.org/abs/1203.2192}.
\end{itemize}
\end{document}
